\SourceLeaf{381}%
auf den Parallelkreis, auf dem der Krümmungsradius des Meridians gleich dem Radius eines Kreises ist, dessen Umfang dem gesamten elliptischen Meridian entspricht. Bei der zugrunde gelegten Abplattung 0,00324 ergibt sich für die Breite dieses Parallelkreises \(45^0\,3'\,29''2\). Die geodätischen Messungen ergeben für den ganzen Bogen dieses Parallelkreises zwischen den Meridianen von Mailand (Sternwarte) und Usson 475121,06 Meter, während die Beobachtungen der Pulversignale für den Längenunterschied \(6^0\,1'\,41''7\) ergeben haben. Diese Zahlen kann man als das Hauptergebnis der Messungen betrachten. Durch den Vergleich eines solchen Längengradbogens mit dem Ergebnis einer Breitengradmessung lässt sich theoretisch die Erdabplattung bestimmen. Aus einem solchen Vergleich ihres Ergebnisses mit dem Bogen von Greenwich bis Formentera finden die Verfasser die Abplattung \(\frac{1}{241}\). Wie wenig zuverlässig dieses Verfahren ist, zeigt sich am auffälligsten, wenn man anstelle des ganzen Bogens die einzelnen Stücke auf ähnliche Weise behandelt. Beim Vergleich desselben Meridianbogens mit dem Stück d’Usson --- Colombier findet der Rezensent so die Abplattung \(\frac{1}{169}\), mit dem zweiten Stück Colombier --- Mont Cenis \(\frac{1}{226}\), mit dem vierten, Turin --- Mailand, \(\frac{1}{110}\), mit dem dritten Stück Mont Cenis --- Turin dagegen eine Verlängerung von \(\frac{1}{82}\). Für die genauere Bestimmung der Erddimensionen ist hiervon also wenig zu erwarten. Umso wichtiger sind die Ergebnisse als neue Bestätigung der Unregelmäßigkeit der Erdgestalt, die gerade in Oberitalien im größten Maßstab auftritt. Am deutlichsten treten diese Unregelmäßigkeiten hervor, wenn man die astronomisch bestimmten Längenunterschiede mit denen vergleicht, die aus den geodätischen Messungen nach einer plausiblen Annahme über die Erdgestalt im Großen berechnet wurden. Die Verfasser haben diese Rechnung mit der Abplattung 0,00324 und dem Äquatorradius 6376986 Meter durchgeführt. So ergeben sich die westlichen Längenunterschiede gegenüber Mailand in Zeitmaß
\begin{center}
\begin{tabular}{l|r@{\quad}l|r@{\quad}l|r}
&\multicolumn{2}{c|}{astronomisch}&\multicolumn{2}{c|}{geodätisch}&Unterschied\\
Turin & \(5'\) & \(58''85\) & \(6'\) & \(0''93\) & \(-2''08\)\\
Mont Cenis & \(9\) & \(0,20\) & \(8\) & \(59,49\) & \(+0,71\)\\
Colombier & \(13\) & \(44,23\) & \(13\) & \(43,84\) & \(+0,39\)\\
D’Usson & \(24\) & \(6,78\) & \(24\) & \(8,02\) & \(-1,24\)
\end{tabular}
\end{center}
